% ch5_b §5.5–5.7 (书页 157–165 / p171–p179)
%--A-TAIL--
%JOIN%
受到修正，而且这种振动模式的晶格频率将依赖于 $\epsilon(\mathbf{q})$（见 §6.11）。于是，$\epsilon$ 中的奇异性将反映在声子频率上。

在固体中传播的任何其他波动系统中，只要它们与传导电子有相互作用，也很可能观察到同样的效应，例如自旋波（§10.10）。原则上，这为研究费米面提供了一种普遍方法——尽管在实践中这一效应未必探测得到。

\section{Friedel 求和规则}

$\epsilon(\mathbf{q},0)$ 中的奇异性还有另一个有趣的后果。假设我们把式 (5.36) 代入点电荷屏蔽势的公式。按照式 (5.18) 与 (5.31)，应当得到

\begin{equation*}
\mathcal{U}(\mathbf{r}) = 4\pi e^2 \int \left\{ q^2 + \frac{4\pi e^2 n}{\frac{2}{3}\mathcal{E}_F} \left[ \frac{1}{2} + \frac{4k_F^2 - q^2}{8k_F q} \ln\left| \frac{2k_F + q}{2k_F - q} \right| \right] \right\}^{-1} e^{-i\mathbf{q}\cdot\mathbf{r}} \,\mathrm{d}\mathbf{q}
\tag{5.38}
\end{equation*}

此即距离 $r$ 处的势。不必真的完成这个计算，我们也很容易相信：$q = 2k_F$ 处的奇异性将作为一项特殊贡献显现在 $\mathcal{U}(\mathbf{r})$ 中；它将不再是光滑的指数函数，而会含有波数为 $2k_F$ 的振荡。

让我们不从式 (5.38) 出发来推导这一效应，而是循着 Friedel 提出的另一条颇为不同的论证路线。考虑一个球对称势 $\mathcal{U}(\mathbf{r})$。我们知道，在这种势存在时，Schrödinger 方程有一个形如下式的解

\begin{equation*}
\psi_{k,l} \sim \frac{1}{r} \sin\left( kr - \tfrac{1}{2}l\pi + \eta_l \right) P_l(\cos\theta)
\tag{5.39}
\end{equation*}

它在远处成立。式中 $P_l(\cos\theta)$ 是 $l$ 阶的普通球谐函数。我们现在考虑的不是散射问题，而是一个“自由”电子的定态；这个空间中除了位于中心的势 $\mathcal{U}(\mathbf{r})$ 之外空无一物。

像 (5.39) 这类解的存在——与 $\mathcal{U}(\mathbf{r}) = 0$ 时的对应解相比，其相位移动了一个量 $\eta_l$——正是求解这类势的散射问题的\emph{分波法}（partial wave method）的核心（参见 §3.9）。众所周知，像 (5.39) 这样的函数可与表示初态的平面波 $\exp(i\mathbf{k}\cdot\mathbf{r})$ 相匹配，而散射部分便可以计算出来。\emph{微分散射截面}由下式给出

\begin{equation*}
\sigma(\theta) = \frac{1}{k^2} \left| \sum_{l=0}^{\infty} (2l+1)\, e^{i\eta_l} \sin\eta_l \, P_l(\cos\theta) \right|^2 .
\tag{5.40}
\end{equation*}

但还是回到我们的定态 (5.39)。假设我们取它们作为体系的基本本征态，并像对待平面波态那样，按能量递增的次序把它们逐个填满。为了便于计数状态，方便的办法是换用另一类边界条件——明显的规则是在 $r = R$ 处 $\psi_k = 0$，就好像我们有一个半径为 $R$ 的大球，原子置于中心。于是 $k$ 的“允许”值必须满足

\begin{equation*}
kR - \tfrac{1}{2}l\pi + \eta_l = \text{integer} \times \pi .
\tag{5.41}
\end{equation*}

倘若我们的球是空的，使得所有相移（phase shift）$\eta_l$ 都为零，那么 $k$ 的“允许”值就应当是集合

\begin{equation*}
k_n = \frac{\left(n + \tfrac{1}{2}l\right)\pi}{R} .
\tag{5.42}
\end{equation*}

当然，对每一个 $n$ 值有许多不同的 $l$ 值，而对每个 $l$ 又有 $(2l+1)$ 个不同的波函数。但是可以证明——事实上，由关于本征值分布的一个普遍定理即可推出——能量小于 $\mathcal{E}_F$（即波矢小于 $k_F$）的态的总数，在这种方案下与 §1.6 中所用的更常见的立方盒情形相同。

%FIG{93}: （译注：原书本图未附图注文字；图中沿 $k$ 轴标出"允许"值，相移使其移动一个量 $\eta(k)/R$，相邻间隔为 $\pi/R$。）

但原子位于中心时，各 $\eta_l$ 不为零，所以 $k$ 的“允许”值并不恰好落在 $k_n$ 处；它们将移动一个量 $\eta_l/R$。而 $\eta_l$ 又随 $k$ 变化。由图 93 容易看出，在 $k$ 轴上的两点 $k$ 与 $k'$ 之间，新的集合将含有

\begin{equation*}
\{\eta_l(k) - \eta_l(k')\}/\pi
\end{equation*}

个新的允许值。

现在假设（这是合理的）$\eta_l$ 在 $k = 0$ 处趋于零。把各个能级填充到同一费米波矢 $k = k_F$ 所需的新电子总数必为

\begin{equation*}
Z = \frac{2}{\pi} \sum_l (2l+1)\, \eta_l(k_F)
\tag{5.43}
\end{equation*}

这里对给定的 $l$ 计入了 $(2l+1)$ 个态，再加上电子的 2 个自旋态。我们已把这个数取作 $Z$，即杂质与溶剂金属之间的价电子数之差，因为恰好需要这么多的额外电子聚集在杂质附近来中和它的电荷。

这就是 \emph{Friedel 求和规则}（Friedel sum rule）。它依据两条普遍原理：杂质的电荷必须在有限距离内被过剩的电子所中和；在远离杂质处，费米波矢 $k_F$ 必须与纯净完整晶体中的相同。

这个公式为计算杂质的散射截面提供了一种极有用的方法。我们选取某个函数，例如带一个可调参数 $\lambda$ 的屏蔽 Coulomb 势 (5.29)。我们算出诸相移，并求出使它们满足 (5.43) 的 $\lambda$ 值。利用这些 $\eta_l$ 值，便可由分波公式 (5.40) 求得散射截面。这一步骤给出相当好的结果，而且对函数 $\mathcal{U}(\mathbf{r})$ 的形状不很敏感。原则上，它肯定比——比如说——在 Born 近似中使用 (5.38) 的矩阵元要好。传导电子太慢，那种做法并不成立。

但现在来考虑与像 (5.39) 那样带相移的波相联系的电子密度。在距原子很远的 $r$ 处，我们将发现多余电荷

\begin{align*}
\delta\rho &= e \sum_l (2l+1) \int_0^{k_F} \left\{ |\psi_{k,l}(\eta_l)|^2 - |\psi_{k,l}(0)|^2 \right\} \frac{2R}{\pi}\, \mathrm{d}k \\
&= \frac{e}{2\pi^2} \sum_l (2l+1) \int_0^{k_F} \left\{ \sin^2\left(kr - \tfrac{1}{2}l\pi + \eta_l\right) - \sin^2\left(kr - \tfrac{1}{2}l\pi\right) \right\} \frac{1}{r^2}\, \mathrm{d}k \\
&= \frac{e}{2\pi^2} \sum_l (2l+1) \int_0^{k_F} \sin\left(2kr - l\pi + \eta_l\right) \sin\eta_l \, \frac{1}{r^3}\, \mathrm{d}(kr) \\
&\sim \frac{e}{2\pi^2} \sum_l (2l+1) \left(-1\right)^l \sin\eta_l \, \frac{\cos(2k_F r + \eta_l)}{r^3}
\tag{5.44}
\end{align*}

这里对半径为 $R$ 的球中的波函数 (5.39) 使用了归一化因子 $(2\pi R)^{-\frac{1}{2}}$，并略去了在 $r$ 很大时更快趋于零的各项。

这就是与 (5.38) 中的奇异性相联系的振荡电荷密度。它仅以 $1/r^3$ 的方式趋于零，因此绝非可忽略的效应。有趣的是，杂质附近的电荷并不是简单地堆积起来；存在 $\delta\rho$ 为负的区域，在那里电子实际上被驱赶出去。这一点在与杂质间相互作用以及合金中 Knight 位移（Knight shift）有关的现象中具有重要意义（见 §10.2）。

实际情况是：这份局域电荷并不是束缚在杂质周围的一两个电子的电荷。它来自费米气体中的快电子被吸引向杂质——它们倾向于在邻近处稍作逗留，并在那里挤得更紧。电子波长的陡峭截断于是引起干涉效应，在散射中心四周投下电荷的晕环。

%FIG{94}: 杂质周围的电荷晕

当\emph{过渡金属}溶入“正常”金属时，我们必须考虑未满的 d 壳层。正如 §3.10 中所见，自由原子的 d 轨道表现为一个\emph{共振}；在其中，如式 (3.81) 那样，\emph{d 相移}在宽度为 $W$ 的能量区间内迅速通过 $\pi/2$。由 Friedel 求和规则 (5.43) 可知，当我们穿越共振时，最多可达 10 个电子的电荷迅速在杂质上积累起来。但杂质上的总电荷必须自行稳定以求中和它的核电荷；离子中“d 电子”的数目必须与自由原子中的大致相同。这意味着溶剂金属的费米能级必须位于共振区之内。

自由原子 d 壳层中的电子总是磁极化的：按照 Hund 规则，它们的自旋排列成与 Pauli 原理相容的最大磁矩。例如在 Mn 中，全部 5 个 d 电子可以取相同的自旋。当原子溶入另一金属时，这种倾向可能延续下来，尽管自旋极化已不能再归因于一个定态原子能级。

例如，假设杂质恰好拥有过剩的“向上”自旋电子。对于另一个自旋相同的电子来说，遇到这个原子时，它看到的是一个位于费米能级之下的共振，并被吸引进去（图 95）。于是“向上”自旋密度得到增强，极化态得以稳定。反之，对“向下”自旋的电子，d--d 交换相互作用把共振的能量推高，从而耗减杂质内向下自旋的密度。净效果必须补偿离子的总电荷；但在合适的溶剂中，我们还可能观察到一种动力学的多电子磁矩，其方向缓慢涨落，而大小与自由原子的磁矩相当。这一过程解释了过渡金属稀合金的许多特殊的电学和磁学性质（见 §10.6）。

%FIG{95}: 磁性杂质处电子的共振，杂质极化沿"向上"方向

\section{半导体的介电常数}

金属的宏观介电常数当 $\mathbf{q} \to 0$ 时趋于无穷，因为金属具有高电导率。这个结果由 (5.16) 和 (5.21) 得出。然而在半导体或绝缘体中，电导率很小；除载流子的热激发之外，它可取为零。静态介电常数 $\epsilon(\mathbf{q},0)$ 在长波下应保持有限。

这源于已满态与空态之间的能隙。形如 (5.16) 的求和中的分母在 $\mathbf{q} = \mathbf{0}$ 处并不为零，而是保持有限。

要计算一个实际半导体的 $\epsilon(\mathbf{q},0)$，需要能带结构的细致模型。我们不能使用 §5.1 的简单理论，因为未微扰的电子态 $|\mathbf{k}\rangle$ 已不再是简单的平面波。不过很容易证明：只需把微扰波 $\exp(i\mathbf{q}\cdot\mathbf{r})$ 的矩阵元的平方引入求和 (5.16)，便得到正确的公式

\begin{equation*}
\epsilon(\mathbf{q},\omega) = 1 + \frac{4\pi e^2}{q^2} \sum_{\mathbf{k},\mathbf{g}} \frac{\left| \langle \mathbf{k}|\, e^{i\mathbf{q}\cdot\mathbf{r}}\, |\mathbf{k}+\mathbf{q}+\mathbf{g}\rangle \right|^2 \left\{ f^0(\mathbf{k}) - f^0(\mathbf{k}+\mathbf{q}+\mathbf{g}) \right\}}{\mathcal{E}(\mathbf{k}+\mathbf{q}+\mathbf{g}) - \mathcal{E}(\mathbf{k}) - \hbar\omega + i\hbar\alpha} .
\tag{5.45}
\end{equation*}

我们必须计入倒格矢 $\mathbf{g} \neq \mathbf{0}$ 的那些项，因为 Bloch 函数之间 $\exp(i\mathbf{q}\cdot\mathbf{r})$ 的矩阵元在二者的波矢相差 $\mathbf{q}+\mathbf{g}$ 时并不为零。这里的情形与 §2.7 的衍射理论本质上相同。

我们不打算精确求值 (5.45)，但可以按下面的办法得到一个粗略近似。让我们采用简约区图式。于是，当 $|\mathbf{k}\rangle$ 是一个已满态时，可用的空态全都在下一个带中；若 $\mathbf{q}$ 很小，$|\mathbf{k}+\mathbf{q}+\mathbf{g}\rangle$ 这时指的是一个几乎竖直地处在 $|\mathbf{k}\rangle$ “正上方”的态。假定能量分母对所有的 $\mathbf{k}$ 值都相同：我们把它取为带隙宽度

\begin{equation*}
\mathcal{E}(\mathbf{k}+\mathbf{q}+\mathbf{g}) - \mathcal{E}(\mathbf{k}) \approx \mathcal{E}_{\mathrm{gap}} .
\tag{5.46}
\end{equation*}

现在我们要对每个 $\mathbf{k}$ 值计算

\begin{equation*}
\sum_{\mathbf{g}} \left| \langle \mathbf{k}|\, e^{i\mathbf{q}\cdot\mathbf{r}}\, |\mathbf{k}+\mathbf{q}+\mathbf{g}\rangle \right|^2
\tag{5.47}
\end{equation*}

为此可以利用下面的定理：——

\begin{equation*}
\sum_n \left( \mathcal{E}_n - \mathcal{E}_s \right) \left| \langle n|\, e^{i\mathbf{q}\cdot\mathbf{r}}\, |s\rangle \right|^2 = \frac{\hbar^2 q^2}{2m},
\tag{5.48}
\end{equation*}

它对 哈密顿量 $\mathcal{H}$ 的本征态 $|n\rangle$ 与 $|s\rangle$ 之间 $\exp(i\mathbf{q}\cdot\mathbf{r})$ 的矩阵元成立，此 哈密顿量的动能项为 $-(\hbar^2/2m)\nabla^2$。这个定理的证明（它是众所周知的\emph{振子强度求和规则}（sum rule for oscillator strengths）的推广）来自把双对易子

\begin{equation*}
\left[ \left[ \mathcal{H},\, e^{i\mathbf{q}\cdot\mathbf{r}} \right],\, e^{-i\mathbf{q}\cdot\mathbf{r}} \right]
\tag{5.49}
\end{equation*}

的期望值在态 $|s\rangle$ 中展开。

如果我们取 $|s\rangle$ 实际上就是我们的态 $|\mathbf{k}\rangle$，并且像 (5.46) 那样，假设对求和中唯一重要的那些态 $|n\rangle$ 有

\begin{equation*}
\mathcal{E}_n - \mathcal{E}_s = \mathcal{E}(\mathbf{k}+\mathbf{q}+\mathbf{g}) - \mathcal{E}(\mathbf{k}) \approx \mathcal{E}_{\mathrm{gap}}
\tag{5.50}
\end{equation*}

那么我们就得到

\begin{equation*}
\sum_{\mathbf{g}} \left| \langle \mathbf{k}|\, e^{i\mathbf{q}\cdot\mathbf{r}}\, |\mathbf{k}+\mathbf{q}+\mathbf{g}\rangle \right|^2 \approx \frac{1}{\mathcal{E}_{\mathrm{gap}}} \, \frac{\hbar^2 q^2}{2m} .
\tag{5.51}
\end{equation*}

把 (5.46) 和 (5.51) 代入 (5.45)，并把满带中 $n$ 个电子各自的贡献加起来（同时也计入 $|\mathbf{k}\rangle$ 为空而 $|\mathbf{k}+\mathbf{q}+\mathbf{g}\rangle$ 为满的那些项），便得

\begin{align*}
\epsilon(\mathbf{q},0) &\approx 1 + \frac{4\pi e^2}{q^2}\, \frac{n}{\mathcal{E}_{\mathrm{gap}}}\, \frac{\hbar^2 q^2}{2m}\, \frac{2}{\mathcal{E}_{\mathrm{gap}}} \\
&= 1 + \frac{4\pi n e^2 \hbar^2}{m\left(\mathcal{E}_{\mathrm{gap}}\right)^2} \\
&= 1 + \left( \frac{\hbar\omega_p}{\mathcal{E}_{\mathrm{gap}}} \right)^2 ,
\tag{5.52}
\end{align*}

其中我们定义了\emph{等离体频率}（plasma frequency）

\begin{equation*}
\omega_p = \left( \frac{4\pi n e^2}{m} \right)^{\frac{1}{2}}
\tag{5.53}
\end{equation*}

——这个参量的物理意义稍后即将显现。

公式 (5.52) 表明，静态介电常数在 $\mathbf{q} = \mathbf{0}$ 处趋于一个常数。这应当就是固体观测到的宏观介电常数。能隙越小，介质越容易极化，即 $\epsilon$ 值越大。当然，我们不能把所取的 $\mathcal{E}_{\mathrm{gap}}$ 值等同于导带底与价带顶之间的光学间隙；它更像是简约区中竖直方向上相互对着的那些态之间的平均能量距离。

于是，半导体中杂质周围的静电场并不像金属中那样受到屏蔽：它只是简单地缩减为 $e^2/\epsilon r$，其中 $\epsilon$ 是静态介电常数。这使得束缚电子态得以在杂质周围形成——这就是\emph{杂质能级}（impurity levels），我们将在 §6.4 中再次讨论。当然，如果介质中存在残余载流子密度 $n_0$，那么在远处我们仍会看到杂质受到经典屏蔽，正如 Debye--Hückel 公式 (5.26)、(5.29) 所给出的那样。

\section{等离体振荡}

现在我们考虑时间相关的效应，其中“外”微扰场的频率 $\omega$ 不为零。这里会出现几个重要现象。最明显的一个是：(5.16) 或 (5.45) 中的 $\hbar\omega$ 恰好与被微扰所混合的一对态之间的能量差相合，也就是说，对占据能级中的某个 $\mathbf{k}$ 值，

\begin{equation*}
\mathcal{E}(\mathbf{k}+\mathbf{q}+\mathbf{g}) - \mathcal{E}(\mathbf{k}) = \hbar\omega .
\tag{5.54}
\end{equation*}

这时 $\epsilon(\mathbf{q},\omega)$ 中出现奇异性——或者说虚部 $i\hbar\alpha$ 变为占优。这恰恰对应于电子的光学吸收（或发射），即电子从态 $|\mathbf{k}\rangle$ 跃迁到态 $|\mathbf{k}+\mathbf{q}+\mathbf{g}\rangle$。我们把这称为体系的一种\emph{单粒子激发}（single-particle excitation）。这些效应我们在 §8.5 中再来讨论。

但假设 $\hbar\omega$ 很大——远大于 (5.16) 分母中的任何一个能量差。我们可以把求和改写成如下形式：

\begin{equation*}
\epsilon(\mathbf{q},\omega) = 1 + \frac{4\pi e^2}{q^2} \sum_{\mathbf{k}} \frac{2f^0(\mathbf{k}) \left\{ \mathcal{E}(\mathbf{k}) - \mathcal{E}(\mathbf{k}+\mathbf{q}) \right\}}{(\hbar\omega)^2 - \left\{ \mathcal{E}(\mathbf{k}) - \mathcal{E}(\mathbf{k}+\mathbf{q}) \right\}^2} .
\tag{5.55}
\end{equation*}

于是可以把分母就取为 $(\hbar\omega)^2$，并把分子按 $\mathbf{q}$ 的幂展开。头几项为零，我们便得到

\begin{align*}
\epsilon(\mathbf{q},\omega) &= 1 + \frac{4\pi e^2}{q^2} \sum_{\mathbf{k}} \frac{f^0(\mathbf{k})}{(\hbar\omega)^2} \left( -q^2 \frac{\partial^2 \mathcal{E}}{\partial k^2} \right) \\
&= 1 - \frac{4\pi n e^2 \hbar^2}{\hbar^2 \omega^2 m} \\
&= 1 - \frac{\omega_p^2}{\omega^2} ,
\tag{5.56}
\end{align*}

其中 $\omega_p$ 是由 (5.53) 定义的等离体频率。

这是对自由电子导出的。但同样的结果也可以由 (5.45) 得到，后者可以改写成

\begin{equation*}
\epsilon(\mathbf{q},\omega) = 1 + \frac{4\pi e^2}{q^2} \sum_{\mathbf{k},\mathbf{g}} \frac{2f^0(\mathbf{k}) \left| \langle \mathbf{k}|\, e^{i\mathbf{q}\cdot\mathbf{r}}\, |\mathbf{k}+\mathbf{q}+\mathbf{g}\rangle \right|^2 \left\{ \mathcal{E}(\mathbf{k}) - \mathcal{E}(\mathbf{k}+\mathbf{q}+\mathbf{g}) \right\}}{(\hbar\omega)^2 - \left\{ \mathcal{E}(\mathbf{k}) - \mathcal{E}(\mathbf{k}+\mathbf{q}+\mathbf{g}) \right\}^2} .
\tag{5.57}
\end{equation*}

如果分母可以当作常数处理，即如果

\begin{equation*}
\hbar\omega \gg \mathcal{E}(\mathbf{k}) - \mathcal{E}(\mathbf{k}+\mathbf{q}+\mathbf{g})
\tag{5.58}
\end{equation*}

对所有被占据的 $\mathbf{k}$，以及对所有矩阵元不可忽略的 $\mathbf{g}$ 都成立，那么我们可以把求和规则 (5.48) 用于对 $\mathbf{g}$ 的求和，对被占据态中 $n$ 个电子的每一个都得到一项 $(4\pi e^2/m\omega^2)$。

公式 (5.56) 表明，当 $\omega \to \omega_p$ 时 $\epsilon \to 0$。而当 $\epsilon = 0$ 时，由 (5.15)，我们有

\begin{equation*}
\mathcal{U} = \frac{\mathcal{V}}{\epsilon} \to \infty .
\tag{5.59}
\end{equation*}

换言之，一个无穷小的“外场” $\mathcal{V}$ 会引起很大的有效场 $\mathcal{U}$：体系是自激发的。频率 $\omega_p$ 是电子气的一个固有振荡模式——称为\emph{等离体模式}（plasma mode）。

这个结果可以用一个初等的论证导出。设体积密度为 $n$ 的电子相对于晶格的固定正电荷背景整体移动一段距离 $\mathbf{x}$。这就建立起一个极化

\begin{equation*}
\mathbf{P} = ne\mathbf{x},
\tag{5.60}
\end{equation*}

它转而又产生一个电场

\begin{equation*}
\mathbf{E} = -4\pi\mathbf{P} .
\tag{5.61}
\end{equation*}

于是每个电荷的运动方程为

\begin{equation*}
m\ddot{\mathbf{x}} = e\mathbf{E} = -4\pi n e^2 \mathbf{x},
\tag{5.62}
\end{equation*}

它定义了频率为 $\omega_p$ 的简谐振荡。

等离体频率对应于相当大的能量，约为 10–20 eV，所以这些模式不会在热能量下被激发。但如果让电子穿过金属薄膜，就可以观察到相应于这些模式中一个或多个量子被激发的能量损失。有趣的是，即使在半导体中，$\omega_p$ 也只依赖于价带中的总电子密度。这是因为只要等离体能量比价带与导带之间的能量差高出相当多——通常后者在 5 eV 以下——条件 (5.58) 就会成立。另一方面，离子芯中的电子不作贡献；它们束缚得太紧。这些能级与导带中的态之间的能量差远大于 $\hbar\omega_p$。

我们可以把 $\epsilon(\mathbf{q},\omega)$ 展开到 $\mathbf{q}$ 的更高次幂，求出使 $\epsilon$ 为零的 $\omega$ 作为 $\mathbf{q}$ 的函数。对自由电子，结果是

\begin{equation*}
\omega(q) = \omega_p \left[ 1 + \frac{3}{10}\, \frac{q^2 v_F^2}{\omega_p^2} + \ldots \right],
\tag{5.63}
\end{equation*}

其中 $v_F$ 是费米面上电子的速度。这样，等离体振荡可以在各种波长下发生——这些模式中的量子可称为\emph{等离激元}（plasmon）。但色散定律 (5.63) 表明，频率对 $q$ 的依赖并不强；等离激元倾向于成为只在介质中缓慢移动的局域化振荡。

严格说来，$\epsilon(\mathbf{q},\omega)$ 是复数：(5.55) 和 (5.57) 表示的是它的实部（假定 $\alpha$ 很小）。虚部由下式给出

\begin{equation*}
\epsilon_2(\mathbf{q},\omega) = \frac{-4\pi^2 e^2}{q^2} \sum_{\mathbf{k}} \left\{ f^0(\mathbf{k}) - f^0(\mathbf{k}+\mathbf{q}) \right\} \delta\!\left\{ \mathcal{E}(\mathbf{k}+\mathbf{q}) - \mathcal{E}(\mathbf{k}) - \hbar\omega \right\},
\tag{5.64}
\end{equation*}

其中的 $\delta$ 函数来自积分 (5.16) 中位于实轴附近的极点

\begin{equation*}
\hbar\omega = \mathcal{E}(\mathbf{k}+\mathbf{q}) - \mathcal{E}(\mathbf{k})
\tag{5.65}
\end{equation*}

处。关于这个公式，有趣的一点是：当

\begin{equation*}
\omega \geqslant q v_F,
\tag{5.66}
\end{equation*}

时它为零；也就是说，此时 $\omega$ 大到仅凭波矢改变 $\mathbf{q}$ 已不能把单个粒子从费米分布中激发出来。这表明等离体模式在这个频率之上不会迅速耗散。

由 (5.56) 还可导出与等离体振荡相联系的另一个有趣现象。如果它表示金属实际的介电常数——如同用长波长的高频电
